\section{绪论}\label{ch:introduction}

\subsection{研究背景及意义}\label{sec:background}

\textbf{（1）多曝光图像融合的应用背景。}

图像融合是将来自同一场景的多幅源图像合成为一幅包含更丰富、更完整信息的融合图像的技术。由于成像设备物理特性的限制，单一图像往往无法同时兼顾亮部和暗部的细节。例如，在多曝光图像融合（Multi-Exposure Image Fusion, MEF）任务中，同一场景在不同曝光条件下拍摄的多幅图像分别捕获了场景的亮部和暗部细节：欠曝光图像保留了亮部区域的纹理信息，但暗部区域细节缺失；过曝光图像则恰好相反，暗部区域细节丰富而亮部区域出现饱和现象。通过融合不同曝光级别的图像，可以生成一幅同时包含高光和阴影细节的高动态范围图像，从而在标准显示设备上呈现更接近人眼真实视觉感知的图像效果。

多曝光图像融合技术在数字摄影、安防监控、医学成像、自动驾驶等众多领域具有广泛的应用价值。在数字摄影领域，MEF技术使得消费级相机能够在高对比度场景下拍摄出细节丰富的照片，克服了传统相机动态范围不足的物理限制。在安防监控领域，MEF技术可以有效解决光照不均场景下的目标检测与识别问题。在医学成像领域，不同模态医学图像（如CT与MRI）的融合可以为临床诊断提供更全面的解剖和功能信息。在自动驾驶领域，红外与可见光图像的融合可以提高复杂光照条件下的环境感知能力。

\textbf{（2）现有方法的局限性。}

近年来，随着深度学习技术的快速发展，基于数据驱动的深度图像融合方法取得了显著的性能提升。卷积神经网络（Convolutional Neural Network, CNN）和Transformer等架构被广泛引入图像融合领域，通过学习源图像的非线性映射关系来实现高质量的融合结果。然而，现有的深度学习方法大多仅依赖于源图像的视觉特征进行融合决策，缺乏对图像语义内容的深层理解。这种纯视觉驱动的融合策略在面对复杂场景时容易产生伪影、对比度失真或语义信息丢失等问题。

\textbf{（3）研究动机与意义。}

视觉语言模型（Vision-Language Model, VLM）的兴起为图像融合领域提供了新的研究思路。大语言模型（Large Language Model, LLM）在理解图像语义内容方面展现出了强大的能力，能够生成准确描述图像关键信息的文本描述。将文本语义信息引入图像融合过程，可以为融合决策提供高层语义指导，使融合结果在保留视觉细节的同时更加符合人类的认知理解。

基于上述背景，本文以FILM（Image Fusion via Vision-Language Model）视觉语言图像融合方法\supercite{zhao2024film}为基础，面向多曝光图像融合任务开展模型复现、推理流程适配和交互式系统实现。该方法通过跨注意力（Cross-Attention）机制将视觉特征与文本特征进行有效对齐和融合。本文在此基础上设计并实现了一个基于Web的交互式图像融合平台，使得用户可以通过直观的界面完成图像上传、融合处理和质量评估等操作流程。本研究的主要意义体现在以下几个方面：

第一，完成了视觉语言图像融合方法在多曝光图像融合任务中的复现与工程适配，验证了文本语义特征对融合细节保持和视觉信息保真的辅助作用。

第二，结合代码实现梳理了基于跨注意力机制的视觉-文本特征交互流程，为后续在图像融合系统中使用多模态语义信息提供了实现参考。

第三，开发了一个功能完整的交互式图像融合系统，将学术研究成果转化为实用的工具平台，促进了多曝光图像融合技术的实际应用推广。

\subsection{国内外研究现状}\label{sec:related_work}

\textbf{（1）传统多曝光图像融合方法。}

早期的多曝光图像融合方法主要基于手工设计的特征和规则。Debevec和Malik\supercite{debevec1997hdr}提出了经典的HDR重建方法，通过多幅不同曝光时间的图像重建场景的辐射度图，然后通过色调映射（Tone Mapping）将高动态范围图像压缩到标准显示范围。该方法提供了一种简单、实用的HDR重建方法，无需特殊硬件，仅需已知曝光时间，逆求解出响应函数，即可将HDR辐射图重构。

基于多尺度分解的方法将源图像分解为不同尺度的子带系数，然后采用特定的融合规则对系数进行合并。小波变换（Wavelet Transform）、拉普拉斯金字塔（Laplacian Pyramid）和对比度金字塔（Contrast Pyramid）等方法被广泛应用于多曝光图像融合。Burt和Adelson\supercite{burt1983laplacian}提出的拉普拉斯金字塔方法对不同曝光度的同一场景图像分别构建拉普拉斯金字塔，然后在每个尺度层上选取或合并最有信息的像素（例如曝光良好的细节），最后用逆变换重建出高动态范围（HDR）或高质量融合图像。

基于稀疏表示的方法假设图像可以由过完备字典中的少量原子线性表示，通过求解稀疏编码来实现图像融合。Li等\supercite{li2011sparse}提出了一种基于稀疏表示的图像融合框架，利用训练得到的过完备字典对源图像进行稀疏分解，然后对稀疏系数进行融合。这类方法在保留图像细节方面表现优异，但计算复杂度较高，难以满足实时应用需求。

基于信息论的方法利用熵、互信息等度量来指导融合过程。Mertens等\supercite{mertens2009exposure}提出了一种无需重建HDR图像的曝光融合方法，可直接从不同曝光的照片序列生成高质量、低动态范围（LDR）的显示图像。该方法计算效率高、实现简单，被广泛应用于消费级相机的HDR模式。

\textbf{（2）基于深度学习的图像融合方法。}

随着深度学习技术的发展，基于CNN的图像融合方法逐渐成为主流。Prabhakar等\supercite{prabhakar2020deep}提出了一种基于深度学习的快速多曝光图像融合方法（MEF-Net），是首个结合全卷积网络与可微分引导滤波的端到端多曝光融合方法。。该方法通过预训练的参考图像进行监督学习，显著提高了融合效率。

Liu等\supercite{liu2020ifcnn}提出了IFCNN（Image Fusion via Convolutional Neural Network），一个通用的深度学习图像融合框架。该方法使用最大L1范数策略融合源图像的卷积特征表示，并通过零填充策略处理任意数量的源图像输入。实验结果表明，IFCNN在多种融合任务（红外-可见光、医学、多曝光）中均取得了优异的性能。

在医学图像融合领域，Zhang等\supercite{zhang2021fusion}提出了一种深度耦合反馈网络（CF-Net），用于联合多曝光融合和图像超分辨率任务。该方法通过耦合反馈机制在不同处理阶段之间共享信息，实现了多任务的协同优化。

Xu等\supercite{xu2020cnn}提出了一种基于CNN特征的多曝光融合方法，利用深度卷积网络提取的多尺度特征进行融合决策。该方法在保留局部细节和全局结构信息方面均表现优异。

\textbf{（3）基于无监督学习的图像融合方法。}

由于获取高质量的参考融合图像成本较高，基于无监督学习的融合方法受到了广泛关注。Li等\supercite{li2022unsupervised}提出了一种基于颜色无参考损失函数的深度无监督多曝光图像融合方法，设计了多种颜色约束的损失函数来指导网络训练，避免了参考图像的依赖。

DPE-MEF（Deep Perceptual Enhancement for Multi-Exposure Fusion）\supercite{xu2022dpe}是一种基于深度感知增强的多曝光融合方法，使用Retinex理论生成参考图像，通过感知损失和像素损失的自动加权组合来训练网络。该方法在保留自然色彩和细节方面取得了优异的性能。

\textbf{（4）视觉语言模型在图像理解中的应用。}

视觉语言模型的快速发展为图像理解任务带来了范式性的变革。CLIP（Contrastive Language-Image Pre-training）\supercite{radford2021clip}通过在大规模图像-文本对上进行对比学习，学会了将图像和文本映射到共享的嵌入空间。BLIP（Bootstrapping Language-Image Pre-training）进一步引入了视觉语言理解和生成能力，能够完成图像描述、视觉问答等多种任务。

BLIP2\supercite{li2023blip2}通过引入Q-Former模块，有效地将预训练的大语言模型（如OPT、FlanT5）与视觉编码器连接起来，在保持大语言模型强大语言能力的同时，实现了高效的视觉语言对齐。BLIP2提取的文本特征具有高度的语义描述能力，可以被用于指导下游的视觉任务。

本文的方法也使用了BLIP2作为文本编码器，将从原图得到的Image Caption、Dense Caption和Segment Anything信息进行有效融合，随后作为文本信息指导过曝和曝光不足的照片进行融合。

\textbf{（5）现有工作分析与总结。}

综合以上文献综述，当前多曝光图像融合领域存在以下主要挑战：

（1）\textbf{语义信息利用不足}：现有的深度学习方法主要依赖于底层视觉特征（如梯度、纹理、对比度）进行融合决策，缺乏对图像语义内容的高层理解。这导致融合结果在某些场景下出现语义不一致的问题，例如在包含人物、建筑等语义重要目标的场景中，纯视觉驱动的融合方法可能无法正确地保留关键语义信息。

（2）\textbf{多模态信息融合效率低}：当引入多种信息源（如图像特征与文本特征）时，如何有效地进行跨模态信息对齐和融合仍然是一个开放性问题。简单的特征拼接或加权求和难以充分挖掘不同模态之间的互补信息。

（3）\textbf{系统实用性有待提升}：大多数研究工作集中在算法层面的性能提升，缺乏对系统级的展示，用户友好性不足，限制了技术的实际应用推广。

针对上述挑战，本文基于FILM视觉语言图像融合方法开展多曝光图像融合任务适配，并实现配套的交互式融合系统，从算法复现、推理部署和系统应用三个层面推动多曝光图像融合技术的落地。

\subsection{本文主要研究内容}\label{sec:content}

本文围绕基于视觉语言模型的多曝光图像融合方法及其应用展开研究，主要研究内容包括以下三个方面：

\textbf{（1）基于视觉语言模型的图像融合方法研究。}

本文围绕文本语义引导的图像融合框架FILM开展复现与任务适配。该方法利用预提取的视觉语言文本语义特征，通过跨注意力机制将文本特征与视觉特征进行对齐和融合。具体而言，本文重点分析和实现以下核心模块：

文本语义特征提取与预处理模块：使用预训练视觉语言模型得到的文本描述特征，通过一维卷积将其投影到与视觉特征兼容的嵌入空间。

视觉-文本跨注意力融合模块：复现基于Restormer Transformer的跨注意力块（restormer\_cablock），将图像特征投影到文本特征空间后，以文本特征作为查询（Query）、图像特征作为键（Key）和值（Value）进行跨注意力计算，实现文本语义对视觉特征的引导作用。

多阶段特征融合模块：采用三级级联的跨注意力块结构，即先进行降采样再进行上采样，逐级深化文本语义引导下的视觉特征融合，最终通过Restormer Transformer块进行特征重建，生成融合亮度图像。

\textbf{（2）多曝光图像融合系统的设计与实现。}

在算法研究的基础上，本文设计并实现了一个基于Web的交互式多曝光图像融合系统。系统采用前后端分离的架构设计，前端基于Vue 3框架构建单页应用（SPA），提供图像上传、算法选择、融合结果可视化、质量评估等交互功能；后端基于FastAPI框架构建RESTful API服务，负责加载预训练的FILM模型并执行融合推理计算。系统支持AI融合（FILM模型）和传统融合算法（均值融合、像素最大值融合、Mertens融合）两种模式，满足不同用户的需求。

\textbf{（3）融合质量评估与实验验证。}

本文设计了一系列定量和定性实验来验证方法适配后的有效性。在定量评估方面，采用信息熵（Entropy, EN）、标准差（Standard Deviation, SD）、空间频率（Spatial Frequency, SF）、平均梯度（Average Gradient, AG）、视觉信息保真度（Visual Information Fidelity, VIF）和边缘信息保持度（Qabf）等六种客观评价指标对融合结果进行评估。对于未在本项目中重新训练的结构变体，本文主要结合已有FILM方法的模块设计进行分析，不将其作为本文独立完成的消融实验结果。

\subsection{本文组织结构}\label{sec:organization}

本文共分为五章，各章内容安排如下：

第一章为绪论。介绍了多曝光图像融合的研究背景和意义，综述了国内外相关研究工作的进展与现状，分析了现有方法的局限性，阐述了本文的主要研究内容和创新点，并给出了全文的组织架构。

第二章为相关理论与技术概述。介绍了深度学习在图像融合领域的应用基础，包括卷积神经网络和Transformer架构的核心概念；详细阐述了视觉语言模型的基本原理及其在图像理解中的应用；介绍了多曝光图像融合任务的特点、YCbCr色彩空间的处理方法以及融合质量评价指标体系。

第三章为基于视觉语言模型的图像融合方法研究。详细阐述了FILM方法的整体架构设计，包括文本特征提取与预处理、视觉-文本跨注意力融合机制、多阶段级联融合结构等核心模块；介绍了损失函数的设计原理和训练策略；给出了实验设置、评估结果和消融分析。

第四章为基于视觉语言模型的图像融合系统设计与实现。介绍了系统的需求分析和整体架构设计，详细描述了前后端的功能模块实现过程，展示了系统的运行效果，并进行了功能验证和性能测试。

第五章为总结与展望。对全文的研究工作进行了系统性总结，概括了主要的创新点和贡献，同时分析了当前研究的不足之处，并对未来的研究方向进行了展望。
